\ficha{Bordo e Kydland (1990), The Gold Standard as a Rule}{bordokydland1990}

\begin{identificacao}
BORDO, Michael D.; KYDLAND, Finn E. The gold standard as a rule: an essay in
exploration. \textit{Explorations in Economic History}, v. 32, n. 4,
p. 423-464, 1995.

Artigo teórico e de história comparada, em inglês. Leitura integral do working
paper que precede a versão publicada: NBER Working Paper n. 3367, Cambridge,
National Bureau of Economic Research, maio de 1990, 49 páginas de texto, duas
tabelas, três figuras e notas finais. As páginas citadas nesta ficha são as do
working paper de 1990, não as do artigo de 1995.
\end{identificacao}

\bloco{Problema e tese}
Por que alguns países sustentaram a conversibilidade em ouro por décadas e
outros a abandonaram sempre que ela custou caro? O padrão-ouro não é tratado
como mecanismo automático de estabilidade de preços nem como simples arranjo
cambial, mas como tecnologia de compromisso que impede o governo de alterar a
política planejada para o futuro. A tese é que a regra seguida antes de 1914
pela Inglaterra e, em menor medida, pelos Estados Unidos era uma regra
contingente: autorizava suspender o preço fixo do ouro em emergência de guerra
sob a condição de restabelecer a conversibilidade à paridade original quando a
emergência passasse. Em outros países, a mesma regra valeu mais como meta
desejável do que como restrição operacional.

\bloco{Argumento}
\begin{enumerate}
\item O conceito de regra é reformulado. A Currency School associava regra a
impessoalidade e automaticidade, com a emissão fiduciária variando com a
reserva de ouro (p. 1). Aqui regra passa a significar vinculação de ações de
política ao longo do tempo, na chave da inconsistência temporal de Kydland e
Prescott (p. 2).

\item O que a regra resolve não é sobretudo a inflação, é a dívida. Sem
compromisso, o governo tem incentivo a tributar o capital já instalado e a dar
calote na dívida, inclusive por inflação, e o público antecipa isso no preço
dos títulos. O padrão-ouro entra como instituição que amarra, análoga à lei de
patentes, com a vantagem de ser simples e transparente (p. 8).

\item A regra é contingente por desenho. O soberano mantém o preço fixo exceto
em guerra maior, quando pode suspender pagamentos em espécie e emitir papel, no
entendimento de que a dívida será paga em ouro depois (p. 12). Restringir a
conversibilidade de depósitos em pânico bancário é coisa distinta de suspender
a conversibilidade externa (p. 12).

\item O teste é narrativa comparada. Inglaterra de 1821 a 1914 é o único caso
de observância contínua; as suspensões de 1797 a 1821 e de 1914 a 1925 operam
a regra contingente, e os adiamentos da retomada depois de 1815 são lidos como
discrição, embora a retomada efetiva mostre que observar a regra prevalecia
(p. 17). Nos Estados Unidos, os greenbacks de 1862 a 1879 também se enquadram,
com apreciação esperada do papel medida pelo diferencial de juros entre títulos
pagos em ouro e em greenbacks (p. 22); a desvalorização de 1933 é ruptura.

\item A assimetria é o resultado que mais interessa ao capítulo. França e
Alemanha se comportam como o centro até 1914. Itália e Argentina seguiram a
regra apenas no nome, a Argentina alternando abandono e retorno conforme o
ciclo de preços de 1880 a 1913, e o restante da América Latina classificado,
com apoio em Fishlow, como caso de prodigalidade monetária e fiscal (p. 28).

\item O enforcement é reputacional e hegemônico. As ``regras do jogo'' foram
violadas com frequência fora da Inglaterra (p. 29) e o Banco da Inglaterra
opera como líder em jogo de Stackelberg (p. 30). O mecanismo que prende a
periferia é outro:
\chave{many developing countries joined the gold standard in order to have
access to the London capital market. Once on the standard, such countries
feared the consequences of departure}{p. 31}
\end{enumerate}

\bloco{Conceitos e definições operacionais}
Padrão-ouro: cada país define o preço do ouro em sua moeda e o mantém fixo,
comprando e vendendo ouro em quantidade ilimitada (p. 11). Padrão misto: no
século XIX vigora a combinação de moeda fiduciária e metálica, e a regra exige
que a fiduciária seja livremente conversível ao preço fixo (p. 11).
Bimetalismo: variante da mesma regra, porque o que importa é o valor fixo da
unidade de conta (p. 11). Regra contingente: compromisso com cláusula de escape
restrita a um estado do mundo bem definido, a guerra, mais prazo de ajuste
conhecido de antemão (p. 12). Discrição:
\chave{Discretion is any purposeful deviation, under whatever guise, from such a
rule}{p. 8}
Credibilidade não é definida, é medida, por diferenciais de juros e erros de
previsão de apreciação (p. 22) e pela entrada de capital privado nas crises
inglesas de 1890 e 1907 (p. 19). Currency board aparece uma única vez, como
forma pela qual colônias se ligavam às moedas metropolitanas (p. 28).

\bloco{Evidência e método}
Não há econometria. O texto encadeia um modelo de tributação ótima com dívida,
de Kydland e Prescott (1980a) e de Lucas e Stokey (1983), e uma resenha
histórica do início do século XIX a 1933, cobrindo Inglaterra, Estados Unidos,
França, Alemanha, Áustria-Hungria, Rússia, Suécia, Itália, Espanha, domínios
britânicos e Argentina. As séries próprias são três: câmbio de Londres sobre
Hamburgo, 1797 a 1821; dólar-libra, 1914 a 1925; preço do ouro em greenbacks,
1862 a 1878. A Tabela 1 decompõe o financiamento das guerras napoleônicas e da
Primeira Guerra na Grã-Bretanha; a Tabela 2 reproduz cálculo de Calomiris. O
resto é literatura secundária.

\bloco{Agentes e vertentes}
Currency School e Banking School (p. 1); bulionistas e antibulionistas em torno
do Relatório do Ouro de 1810 e do Peel's Act de 1819; Bagehot e a doutrina da
responsabilidade; Keynes contra a retomada de 1925 à paridade antiga; nos
Estados Unidos, hard money contra soft money, greenbackers, populistas e a
campanha pela prata livre. Entre os comentadores, Eichengreen para a hegemonia
inglesa, Ford para a Argentina, Fishlow para a América Latina, Fratianni e
Spinelli para a Itália, North e Weingast para a origem do compromisso inglês.

\bloco{Críticas e limites}
Os autores registram que a resenha não explica de forma conclusiva por que a
regra durou em uns países e não em outros, e apresentam estabilidade política e
centralidade inglesa como conjecturas (p. 32). O modelo da seção II não é
calibrado nem estimado, e a passagem para a narrativa é interpretativa: adesão
vira evidência de credibilidade e credibilidade explica adesão. O mecanismo
pelo qual países não centrais aderem ao ouro é asserido, não medido, o que a
literatura posterior foi buscar em prêmios de risco soberano. O Brasil não é
mencionado uma única vez. A explicação por interesses e distribuição de renda,
com a sugestão de que o sufrágio restrito freava as forças inflacionistas, fica
confinada à nota final 38. E a chave normativa que sustenta a comparação,
\chave{Abandoning gold convertibility was viewed as a serious breach of
contract}{p. 32}
mede a periferia pela vara inglesa, sem perguntar se a cláusula de contingência
podia ter conteúdo distinto onde nunca houve conversibilidade estável a
restaurar.

\begin{dialogo}
\item Vocabulário do compromisso em 1906: ``fé pública'', ``contrato'',
``crédito no exterior'', ``confiança do capital estrangeiro''. Defesa da Caixa
apoiada no acesso ao mercado de Londres confirma o mecanismo que os autores
atribuem à periferia (p. 31); defesa apoiada na renda cafeeira e no nível da
taxa favorável à exportação indica motivo distinto, e a hipótese não transfere.
\item Amarra ou instrumento: nos debates sobre o limite de emissão, os fundos
de garantia e de resgate e a competência da Caixa frente ao Tesouro, verificar
se ela é apresentada como restrição que retira do Executivo a faculdade de
emitir. A alteração da taxa em 1910, para 16 pence, é o caso decisivo: defendê-la
por conveniência conjuntural é discrição no sentido do texto (p. 8).
\item Suspensão de 1914: verificar se os jornais tratam o fim do troco como
suspensão temporária com promessa de retorno à mesma taxa, o que aproximaria o
caso brasileiro da cláusula de guerra (p. 12), ou como fim do experimento.
\item Comparação argentina: o texto põe a Argentina entre os que seguiram a
regra só no nome (p. 28 e p. 31). Levantar se a imprensa invoca a Caja de
Conversión como modelo bem-sucedido ou como advertência.
\item Rothschild, funding loan de 1898 e crédito em Londres nos argumentos
pró-Caixa: testa se o medo das consequências da saída do ouro (p. 31) aparece
explicitado como razão para manter a conversibilidade.
\end{dialogo}

\usonocapitulo{Parte I, como fundação conceitual. Sustenta a afirmação de que
``padrão-ouro'', no debate de 1906, não designa apenas um regime cambial, mas um
compromisso institucional sobre a dívida e a emissão, e fornece o par regra e
discrição com que se lê a controvérsia sobre o limite de emissão e sobre a
alteração da taxa. Introduz a assimetria centro-periferia que Bordo e Rockoff,
Obstfeld e Taylor e Ferguson e Schularick detalham na mesma parte. O texto não
trata do Brasil em momento algum e não pode ser citado como evidência sobre o
caso brasileiro; entra como moldura e como contraponto, porque a categoria de
regra contingente foi derivada de países centrais com conversibilidade prévia,
e a Caixa de Conversão converte a uma taxa muito abaixo do par legal de 27
pence, isto é, institui compromisso novo em vez de restaurar um antigo.}
